\begin{abstract}
We simulate a BTW sandpile cellular automaton as a model for self-organized
criticality. We study dimensions $d=2,3,4$ with open and closed boundary
conditions and with conservative and nonconservative driving. For each
configuration, the system is driven into the stationary critical state using
a burn-in procedure monitored by the mean slope $\langle z\rangle(\tau)$.
We then record a large number of avalanches and measure the avalanche
lifetime $t$, size $s$, and linear extension $l$. From these data we estimate
the probability density functions and conditional expectation values, and we
extract the scaling exponents $\alpha$, $\tau$, $\lambda$ and $\gamma_i$ from
weighted linear fits on double-logarithmic plots in the scaling region.
All configurations show approximate power-law behavior over an intermediate
range with finite-size cutoffs for large events. The measured exponents
depend on the driving scheme, while open and closed boundaries give similar
results in the nonconservative case. Our exponent estimates are broadly
consistent with the values reported by Christensen, Fogedby, and Jensen,
with the largest uncertainties and deviations in $d=4$ where the scaling
range is shortest and statistics are more limited.
\end{abstract}